% from ega1-1-fr.tex
\oldpage[I]{95}
canonique fonctoriel
$(B_{[\varphi]})^{\sim}\xrightarrow{\sim}\Phi_*(\widetilde{B})$ est un
isomorphisme de $\mathscr{O}_{X'}$-Algèbres~; si $M$ est un $B$-module,
l'isomorphisme canonique fonctoriel
$(M_{[\varphi]})^{\sim}\xrightarrow{\sim}\Phi_*(\widetilde{M})$ est un
isomorphisme de $\Phi_*(\widetilde{B})$-Modules.

\begin{corollary}[1.6.4]
\label{I.1.6.4-fr}
Le foncteur image directe $\Phi_*$ est exact sur la catégorie des
$\mathscr{O}_X$-Modules quasi-cohérents.
\end{corollary}

En effet, il est clair que $M_{[\varphi]}$ est un foncteur exact en $M$ et
$\widetilde{M'}$ un foncteur exact en $M'$ (\hyperref[I.1.3.5-fr]{1.3.5}).

\begin{proposition}[1.6.5]
\label{I.1.6.5-fr}
Soient $N'$ un $A'$-module, $N$ le $A$-module
$N'\otimes_{A'}A_{[\varphi]}$~; il existe un isomorphisme canonique
fonctoriel du $\mathscr{O}_X$-Module $\Phi^*(\widetilde{N'})$ sur
$\widetilde{N}$.
\end{proposition}

Remarquons d'abord que $j:z'\mapsto z'\otimes 1$ est un
$A'$-homomorphisme de $N'$ dans $N_{[\varphi]}$~; en effet, par définition,
pour $f'\in A'$, on a
$(f'z')\otimes 1=z'\otimes\varphi(f')=\varphi(f')(z'\otimes 1)$. On en
déduit (\hyperref[I.1.3.8-fr]{1.3.8}) un homomorphisme
$\widetilde{j}:\widetilde{N'}\to(N_{[\varphi]})^{\sim}$ de
$\mathscr{O}_{X'}$-Modules, et en vertu de
(\hyperref[I.1.6.3-fr]{1.6.3}), on peut considérer que $\widetilde{j}$
applique $\widetilde{N'}$ dans $\Phi_*(\widetilde{N})$. Il correspond
canoniquement à cet homomorphisme $\widetilde{j}$ un homomorphisme
$h=\widetilde{j}^{\#}$ de $\Phi^*(\widetilde{N'})$ dans $\widetilde{N}$
(\hyperref[0.4.4.3-fr]{0, 4.4.3})~; on va voir que pour chaque fibre,
$h_x$ est \emph{bijectif}. Posons $x'={}^a\varphi(x)$ et soit $f'\in A'$
tel que $x'\in D(f')$~; soit $f=\varphi(f')$. L'anneau
$\Gamma(D(f),\widetilde{A})$ s'identifie à $A_f$, les modules
$\Gamma(D(f),\widetilde{N})$ et $\Gamma(D(f'),\widetilde{N'})$ à $N_f$ et
$N'_{f'}$ respectivement~; soient
$s'\in\Gamma(D(f'),\widetilde{N'})$, identifiée à $n'/{f'}^p$ ($n'\in N'$),
$s$ son image par $\widetilde{j}$ dans $\Gamma(D(f),\widetilde{N})$~; $s$
est identifiée à $(n'\otimes 1)/f^p$. Soit d'autre part
$t\in\Gamma(D(f),\widetilde{A})$, identifiée à $g/f^q$ ($g\in A$)~; alors,
par définition, on a
$h_x(s'_{x'}\otimes t_x)=t_x\cdot s_x$
(\hyperref[0.4.4.3-fr]{0, 4.4.3}). Mais on peut identifier canoniquement
$N_f$ à
$N'_{f'}\otimes_{A'_{f'}}(A_f)_{[\varphi_{f'}]}$
(\hyperref[0.1.5.4-fr]{0, 1.5.4})~; $s$ correspond alors à l'élément
$(n'/{f'}^p)\otimes 1$, et la section $y\mapsto t_y\cdot s_y$ à
$(n'/{f'}^p)\otimes(g/f^q)$. Les diagrammes de compatibilité de
(\hyperref[0.1.5.6-fr]{0, 1.5.6}) montrent que $h_x$ n'est autre que
l'isomorphisme canonique
\[
  \label{I.1.6.5.1-fr}
  N'_{x'}\otimes_{A'_{x'}}(A_x)_{[\varphi_x]}
  \xrightarrow{\sim}N_x=(N'\otimes_{A'}A_{[\varphi]})_x.
  \tag{1.6.5.1}
\]

En outre, soit $v:N'_1\to N'_2$ un homomorphisme de $A'$-modules~; comme
$\widetilde{v}_{x'}=v_{x'}$ pour tout $x'\in X'$, il résulte aussitôt de ce
qui précède que $\Phi^*(\widetilde{v})$ s'identifie canoniquement à
$(v\otimes 1)^{\sim}$, ce qui achève la démonstration de
(\hyperref[I.1.6.5-fr]{1.6.5}).

Si $B'$ est une $A'$-algèbre, l'isomorphisme canonique de
$\Phi^*(\widetilde{B'})$ sur
$(B'\otimes_{A'}A_{[\varphi]})^{\sim}$ est un isomorphisme de
$\mathscr{O}_X$-Algèbres~; si en outre $N'$ est un $B'$-module,
l'isomorphisme canonique de $\Phi^*(\widetilde{N'})$ sur
$(N'\otimes_{A'}A_{[\varphi]})^{\sim}$ est un isomorphisme de
$\Phi^*(\widetilde{B'})$-Modules.

\begin{corollary}[1.6.6]
\label{I.1.6.6-fr}
Les sections de $\Phi^*(\widetilde{N'})$ images canoniques des sections $s'$,
où $s'$ parcourt le $A'$-module $\Gamma(\widetilde{N'})$, engendrent le
$A$-module $\Gamma(\Phi^*(\widetilde{N'}))$.
\end{corollary}

En effet, ces images s'identifient aux éléments $z'\otimes 1$ de $N$,
lorsqu'on identifie $N'$ et $N$ à $\Gamma(\widetilde{N'})$ et
$\Gamma(\widetilde{N})$ respectivement (\hyperref[I.1.3.7-fr]{1.3.7}) et
que $z'$ parcourt $N'$.

\begin{env}[1.6.7]
\label{I.1.6.7-fr}
Dans la démonstration de (\hyperref[I.1.6.5-fr]{1.6.5}), on a prouvé en
passant que l'application canonique
(\hyperref[0.4.4.3.2-fr]{0, 4.4.3.2})
$\rho:\widetilde{N'}\to\Phi_*(\Phi^*(\widetilde{N'}))$ n'est autre que
l'homomorphisme $\widetilde{j}$,
